\section{Линейные оболочки. Теорема о размерности линейной оболочки. Теорема о неполном базисе}

\begin{definition}
Пусть $a_1, a_2, \dots, a_k$ --- система векторов линейного пространства $V$. Множество всех возможных линейных комбинаций этих векторов называется \textbf{линейной оболочкой} (или линейным натяжением) данной системы и обозначается $L(a_1, a_2, \dots, a_k)$:
\[
L(a_1, \dots, a_k) = \{ x \in V \mid x = \lambda_1 a_1 + \lambda_2 a_2 + \dots + \lambda_k a_k, \; \lambda_i \in F \}.
\]
Векторы $a_1, \dots, a_k$ называются \textbf{порождающими} для этой оболочки.
\end{definition}

\begin{theorem}
Линейная оболочка $L(a_1, \dots, a_k)$ является подпространством пространства $V$. Базисом этого подпространства служит любая максимальная линейно независимая (ЛНЗ) подсистема системы векторов $\{a_1, \dots, a_k\}$. Следовательно, размерность линейной оболочки равна рангу системы порождающих векторов:
\[
\dim L(a_1, \dots, a_k) = \operatorname{rang}\{a_1, \dots, a_k\}.
\]
\end{theorem}

\begin{proof}
\begin{enumerate}
    \item \textit{Замкнутость:} Пусть $x, y \in L(a_1, \dots, a_k)$ и $\alpha, \beta \in F$. Тогда $x = \sum_{i=1}^k \lambda_i a_i$ и $y = \sum_{i=1}^k \mu_i a_i$. Их линейная комбинация:
    \[
    \alpha x + \beta y = \sum_{i=1}^k (\alpha \lambda_i + \beta \mu_i) a_i \in L(a_1, \dots, a_k).
    \]
    Значит, $L$ --- подпространство.
    \item \textit{Базис:} Пусть $a_{i_1}, \dots, a_{i_r}$ --- максимальная ЛНЗ подсистема (где $r = \operatorname{rang}\{a_1, \dots, a_k\}$). По свойству максимальности, любой другой вектор $a_j$ из исходной системы линейно выражается через $a_{i_1}, \dots, a_{i_r}$. Следовательно, любая линейная комбинация всех $k$ векторов может быть переписана как линейная комбинация только этих $r$ векторов. Таким образом, они порождают $L$ и по определению образуют в нем базис.
\end{enumerate}
\hfill$\blacksquare$
\end{proof}

\begin{theorem}[О неполном базисе]
Пусть $\dim V = n$, и система векторов $e_1, e_2, \dots, e_k$ ($k < n$) линейно независима в $V$. Тогда эту систему можно дополнить до базиса всего пространства $V$, то есть существуют векторы $e_{k+1}, \dots, e_n \in V$ такие, что система $e_1, \dots, e_k, e_{k+1}, \dots, e_n$ образует базис $V$.
\end{theorem}

\begin{proof}
Рассмотрим подпространство $L = L(e_1, \dots, e_k)$. Так как векторы ЛНЗ, $\dim L = k$. Поскольку $k < n$, то $L \neq V$. 
Следовательно, существует вектор $e_{k+1} \in V$, не принадлежащий $L$. По теореме о добавлении вектора, система $e_1, \dots, e_k, e_{k+1}$ будет линейно независимой. 
Если $k+1 = n$, процесс завершен. Если $k+1 < n$, повторяем рассуждение для подпространства $L(e_1, \dots, e_{k+1})$. Процесс обязательно завершится за конечное число шагов, когда количество векторов достигнет $n$, так как в $n$-мерном пространстве любая система из $n$ ЛНЗ векторов образует базис.
\hfill$\blacksquare$
\end{proof}
